\section{策略安全与系统化部署}

\subsection{Risk-aware fine-tuning}

\begin{figure}[H]
\centering
\includegraphics[width=0.88\textwidth]{slide-30.jpg}
\caption{Lecture slide 30 总结的风险感知 pipeline：先用 conservative critic 约束值，再在部署前做 safety check。}
\end{figure}

课程提到一个安全层：在 actor update 之前检查 critic 输出，如果当前策略值突破 pre-defined threshold 就回滚上一步。这类似 trust region 的 analog，但更偏向 deployment guard。

\begin{warningbox}{Conservative critic 的建议}
保存多组 critic checkpoints，部署时选取 conservative 一组。这样即便 newest critic overestimates，actor 也不会被过强信号误导。
\end{warningbox}

\subsection{在线监控与 rollback}

在生产系统中，常见做法是把 KL、critic loss 和 entropy 平均值推到 dashboard。若任一指标连续 3 次超阈，就自动触发 rollback 和 entropy bonus 重启。

\begin{knowledgebox}{生产系统的经验}
LLM 或 RL 系统里，policy divergence 的 early warning 信号非常重要。课程建议用 3-day rolling window 统计 KL，如果超过 2x historic std，就暂停 actor 更新，让 critic catch up。
\end{knowledgebox}

\subsection{本章小结}

安全部署要求 actor-critic 系统在更新与 inference 之间形成防火墙：conservative critic，guard rails，自动 rollback 与持续监控。

